% Part: counterfactuals
% Chapter: introduction
% Section: strict-conditional

\documentclass[../../../include/open-logic-section]{subfiles}

\begin{document}

\olfileid[zh]{cnt}{int}{str}

\olsection{严格条件句}

Lewis 引入了\emph{严格条件句}~$\strictif$，并论证说，与实质条件句相对的是它——而非实质条件句——才对应蕴含。在真性模态逻辑中，$!A \strictif !B$ 可定义为 $\Box(!A \lif
!B)$。因此，严格条件句（在世界处）为真当且仅当相应的实质条件句是必然的。

严格条件句面对实质条件句的种种悖论表现如何？有着假前件的严格条件句和有着真后件的严格条件句，可能为真，也可能为假。此外，$(!A \strictif !B) \lor (!B \strictif !A)$ 并不有效。严格条件句 $!A \strictif !B$ 也不等价于 $\lnot !A \lor !B$，因此它不是真值函项的。

我们有：
\begin{align}
  !A \strictif !B & \Entails \lnot !A \lor !B
  \text{ 但：}\\
  \lnot !A \lor !B & \Entails/ !A \strictif !B\\
  !B & \Entails/ !A \strictif !B\\
  \lnot !A & \Entails/ !A \strictif !B\\
  \lnot(!A \lif !B) & \Entails/ !A \land \lnot !B
  \text{ 但：}\\
  !A \land \lnot !B & \Entails \lnot(!A \strictif !B)
  \intertext{然而，严格条件句仍支持分离规则：}
  !A, !A \strictif !B & \Entails !B
  \intertext{严格条件句是合聚的：}
  !A \strictif !B,  !A \strictif !C & \Entails !A \strictif (!B \land !C)
  \intertext{前件强化对严格条件句成立：}
  !A \strictif !B & \Entails (!A \land !C) \strictif !B
  \intertext{严格条件句也是传递的：}
  !A \strictif !B, !B \strictif !C & \Entails !A \strictif !C
  \intertext{最后，严格条件句等价于其逆否形式：}
  !A \strictif !B & \Entails \lnot !B \strictif \lnot !A\\
  \lnot !B \strictif \lnot !A & \Entails !A \strictif !B
\end{align}

\begin{prob}
  给出对严格条件句不成立的衍生关系的 \Log{S5}-反例，即对如下各项：
  \begin{enumerate}
  \item $\lnot p \Entails/ \Box(p \lif q)$
  \item $q \Entails/ \Box(p \lif q)$
  \item $\lnot\Box(p \lif q) \Entails/ p \land \lnot q$
  \item $\Entails/ \Box(p \lif q) \lor \Box(q \lif p)$
  \end{enumerate}
\end{prob}

\begin{prob}
  通过给出如下各项的 \Log{S5}-证明来表明严格条件句成立的有效衍生关系：
  \begin{enumerate}
  \item  $\Box(!A \lif !B) \Entails \lnot !A \lor !B$
  \item $!A \land \lnot !B \Entails \lnot\Box(!A \lif !B)$
  \item $!A, \Box(!A \lif !B) \Entails !B$
  \item $\Box(!A \lif !B), \Box(!A \lif !C) \Entails \Box(!A \lif (!B
    \land !C))$
  \item $\Box(!A \lif !B) \Entails \Box((!A \land !C) \lif !B)$
  \item $\Box(!A \lif !B), \Box(!B \lif !C) \Entails \Box(!A
    \lif !C)$
  \item $\Box(!A \lif !B) \Entails \Box(\lnot !B \lif \lnot !A)$
  \item $\Box(\lnot !B \lif \lnot !A) \Entails \Box(!A \lif !B)$
  \end{enumerate}
  \end{prob}

然而，严格条件句仍有着它自己的「悖论」。正如有着假前件或真后件的实质条件句为真一样，有着\emph{必然}假前件或必然真后件的严格条件句也为真。此外，任何真的严格条件句都是必然真，任何假的严格条件句都是必然假。换言之，我们有
\begin{align}
  \Box \lnot !A & \Entails !A \strictif !B\\
  \Box !B & \Entails !A \strictif !B\\
  !A \strictif !B& \Entails \Box(!A \strictif !B)\\
  \lnot(!A \strictif !B) & \Entails \Box\lnot(!A \strictif !B)
\end{align}
如果你把 $\strictif$ 想作「蕴含」，这些便不是问题。逻辑衍生关系终归是数学事实，因此不能是偶然的。但如果你想用 $\strictif$ 作为一个旨在刻画「如果……那么……」的逻辑联结词，它们确实会带来麻烦，尤其是最后两条。因为当然存在着偶然真或偶然假的「如果……那么……」语句——事实上，它们通常既非必然也非不可能。

\begin{prob}
  给出如下各项在 \Log{S5} 中的证明：
  \begin{enumerate}
    \item  $\Box \lnot !A \Entails !A \strictif !B$
    \item $!A \strictif !B \Entails \Box(!A \strictif !B)$
    \item $\lnot(!A \strictif !B)  \Entails \Box\lnot(!A \strictif !B)$
  \end{enumerate}
  利用 $\strictif$ 的定义来完成。
\end{prob}

\end{document}
